\begin{frame}{Cuántas colisiones}
  Modelamos cada hash como un número uniforme entre \(0\) y \(2^{64}-1\).
  Dos elementos distintos colisionan con probabilidad \(2^{-64}\).

  \medskip
  Sea \(X\) la cantidad de colisiones entre \(N\) elementos.
  Por linealidad de la esperanza:
  \[
    \mathbb{E}[X] = \binom{N}{2}\, 2^{-64}.
  \]
  \(X\) es un entero no negativo, así que \(\mathrm{P}(X \neq 0) \le \mathbb{E}[X]\).

  \medskip
  Para \(N = 10^{6}\),
  \[
    \mathrm{P}(X \neq 0)
      \le \binom{10^{6}}{2}\, 2^{-64}
      \approx 2.71 \times 10^{-8}.
  \]
\end{frame}

\begin{frame}{Bases y módulos}
  Para una base y un módulo fijos se pueden armar strings con muchas colisiones.
  Ese caso de prueba tiene que estar construido para esa base.

  \medskip
  En Latam eso casi no aparece.
  Una base grande, elegida al azar, corta los tests que necesitan conocer \(b\).

  \medskip
  Thue-Morse colisiona módulo \(M\) cuando \(b^{2^{k}} \equiv 1 \pmod{M}\) para algún \(k\).
  El módulo no tiene que ser una potencia de \(2\).
  Si es \(2^{64}\) y la base es impar, colisiona igual.
  La vemos al final.
\end{frame}
